\chapter{System Architecture}
\label{ch:architecture}

This chapter describes the high-level architecture of the proposed
driver stress detection and adaptive cabin control system. The
objective is to give the reader a coherent mental model of how
biosignals captured at the wrist are transformed, in fewer than
two seconds of wall-clock time, into a cabin lighting, ventilation,
and audio response inside the vehicle. Sections~\ref{sec:arch-overview}
and \ref{sec:arch-pipeline} describe the system from two complementary
viewpoints: the component view (who talks to whom) and the pipeline
view (what the data looks like at each stage). Section~\ref{sec:arch-hw}
anchors these abstractions to the concrete hardware used in the
prototype, and Section~\ref{sec:arch-compare} positions the architecture
against the prior AlcoWatch BAC system~\citep{acn1408}.

\section{Overall Architecture}
\label{sec:arch-overview}

The system decomposes naturally into five cooperating components,
arranged in a strict data-flow pipeline from sensors to actuators.
Figure~\ref{fig:system-overview} shows the physical layout, the
logical data flow, and the software/hardware boundary between the
wearable and the vehicle module.

\begin{figure}[H]
\centering
\includegraphics[width=\linewidth]{system_overview}
\caption{End-to-end system overview. A wrist-worn Wear~OS device
collects PPG, EDA, skin temperature, and three-axis acceleration,
applies a bidirectional-LSTM + attention classifier on-device, and
transmits a twelve-byte BLE GATT update every thirty seconds. An
Arduino Nano 33 BLE inside the vehicle subscribes to the
characteristic and drives four cabin profiles (Calm, Mild, Moderate,
Critical) via an RGB LED strip, a PWM-controlled fan, and an audio
output.}
\label{fig:system-overview}
\end{figure}

The five components are, from the wrist outward:

\begin{enumerate}
  \item \textbf{Sensor subsystem.} A Wear OS device exposes the four
  physiological channels of interest through the Android \texttt{SensorManager}
  and its health-services SDK. Each channel is sampled at the highest
  rate the platform supports (typically 64~Hz for BVP) and then
  resampled to a common four-hertz rate in the application layer.

  \item \textbf{Classifier subsystem.} A TensorFlow Lite interpreter
  \citep{tflite2023} wraps the bidirectional-LSTM + attention model
  trained offline. It consumes windows of shape
  $[1, 30, 5]$ and emits a four-dimensional softmax distribution
  every half-second internally; the subsystem exposes a running
  average to the BLE layer.

  \item \textbf{BLE peripheral.} A GATT server exposes one primary
  service (\texttt{ABCD1234-1234-5678-1234-56789abcdef0}) containing a
  single notifiable characteristic
  (\texttt{ABCD1235-1234-5678-1234-56789abcdef0}) that carries a
  twelve-byte stress packet. The characteristic is updated every
  thirty seconds.

  \item \textbf{BLE central on the Arduino Nano 33 BLE.} This
  microcontroller scans for the AlcoWatch Stress service, subscribes
  to notifications, parses the twelve-byte packet, and updates its
  internal finite-state machine.

  \item \textbf{Cabin actuator subsystem.} The Arduino drives an
  RGB LED strip (PWM, three colour channels), a cooling fan (PWM
  duty cycle, four discrete levels), and an audio output (tone
  pattern) through three independent general-purpose output paths.
  It falls back to a safe \texttt{WAITING} state if no BLE packet
  is received within sixty seconds.
\end{enumerate}

Two properties of this decomposition deserve emphasis. First,
\emph{inference is fully on-device}: no physiological data ever
leaves the wrist, and the wearable operates correctly even when the
vehicle has no cellular connectivity. Second, the interface between
wrist and vehicle is a \emph{narrow, documented, twelve-byte
characteristic}, which makes the two halves independently replaceable.
In principle, any BLE-capable wearable that emits the same packet
layout can drive the Arduino, and any BLE-capable vehicle module
that consumes the same characteristic can be driven by the wearable.

\section{Data Flow Pipeline}
\label{sec:arch-pipeline}

We now trace a single ``packet'' of data through the system to
describe the signal processing chain end-to-end.

At $t=0$~ms the Wear~OS application receives a new PPG sample from
the hardware sensor at its native 64~Hz rate. The application
accumulates samples in a ring buffer. Every 250~ms (i.e.\ four times
per second) the PPG buffer is polyphase-downsampled to a single
4~Hz sample, peak-detected to update the running inter-beat-interval
(IBI) estimate, and appended to the 30-sample sliding window
together with the latest EDA, skin temperature, and accelerometer
magnitude readings. The window has 50\% overlap: it advances by
fifteen samples (3.75~seconds of wall-clock time) at every inference
step. This choice trades off responsiveness against label stability.

Once per second -- or more precisely, every fifteen samples -- the
full $30 \times 5$ window is z-score normalised using the training-set
per-feature means and standard deviations (cached as constants in
the application) and passed to the TensorFlow Lite interpreter.
The interpreter invokes the bidirectional-LSTM + attention model and
returns a softmax of length four. The argmax of this distribution is
the predicted class, and the corresponding probability is the
confidence.

These per-second predictions are not emitted directly on the BLE
characteristic. Instead, they are accumulated in a thirty-sample
ring, and a majority-vote smoothing step yields a single stress
class and an averaged confidence every thirty seconds. This
temporal smoothing is the key mechanism that keeps the cabin from
flickering between profiles at every short-lived motion artefact,
and it is calibrated to the timescale over which autonomic responses
are known to stabilise.

The resulting decision is packed into the twelve-byte BLE payload
(eight bytes of little-endian timestamp, one byte of stress level,
one byte of confidence rounded to a percentage, two bytes reserved
for future use) and pushed to the Arduino as a GATT notification.
On the Arduino side, the firmware reads the packet, resets its
internal BLE-timeout counter, transitions its finite-state machine
into the new state, and writes the appropriate PWM duty cycles to
the LED, fan, and audio pins. Total latency from the start of the
window to the cabin reacting is dominated by the thirty-second
emission cadence itself; the inference and BLE components contribute
fewer than one hundred milliseconds combined.

\section{Hardware Platform}
\label{sec:arch-hw}

The prototype targets two off-the-shelf hardware platforms, one
on each side of the BLE link.

\textbf{Wearable side.} The reference Wear~OS device is the Samsung
Galaxy Watch class of hardware (Exynos W-series system-on-chip,
1.5~GB RAM, Wear~OS~4+), though any Wear~OS device with minimum
API level 30 that exposes the PPG, skin-temperature, EDA (or a
derivable surrogate), and three-axis accelerometer through the
standard \texttt{Sensor} API will function. The application occupies
roughly 22~MB of storage including the bundled 72~KB TensorFlow Lite
model. Power budget measurements are outside the scope of this
project, but the thirty-second emission cadence and the small model
size have been selected specifically to keep the application well
inside the budget available on a watch that is already running the
user's normal workload.

\textbf{Vehicle side.} The reference vehicle-side platform is the
Arduino Nano 33 BLE. This board is built around the Nordic nRF52840
system-on-chip and exposes an integrated BLE 5 radio, twenty-two
digital I/O pins, eight analogue inputs, and roughly 248~KB of
program memory and 32~KB of SRAM. The firmware developed for this
project occupies approximately 60~KB of program space and less than
8~KB of SRAM, leaving ample headroom for future extensions such as
CAN-bus interfacing or additional sensor inputs. The board is
powered from the vehicle's accessory circuit via a 12~V-to-5~V
regulator; it draws less than 50~mA in its highest-power state.

All other components in the cabin prototype -- the RGB LED strip
(common-anode, PWM-driven through three 220~$\Omega$ current-limiting
resistors), the 5~V axial fan (PWM-controlled through a logic-level
MOSFET), and the piezoelectric tone emitter -- are standard hobbyist
components. No bespoke hardware was fabricated for the project.

\section{Comparison with AlcoWatch BAC Detection System}
\label{sec:arch-compare}

Although this project targets a different physiological phenomenon
(stress rather than alcohol intoxication) and a different intervention
profile (graduated comfort rather than binary interlock), it builds
architecturally on the prior AlcoWatch system filed under provisional
patent ACN1408 \citep{acn1408}. Table~\ref{tab:alcowatch-compare}
contrasts the two systems directly.

\begin{table}[H]
\centering
\caption{Architectural comparison between the prior AlcoWatch BAC
system \citep{acn1408} and the present driver-stress system.}
\label{tab:alcowatch-compare}
\small
\begin{tabular}{p{4cm}p{4cm}p{4cm}}
\toprule
\textbf{Aspect} & \textbf{AlcoWatch BAC \citep{acn1408}} &
\textbf{This Work (Stress)} \\
\midrule
Sensors used & PPG, EDA, skin temperature, ambient temp, humidity
(6 features) & PPG (BVP+IBI), EDA, skin temperature, ACC magnitude
(5 features) \\
ML model type & BiLSTM + attention, regression head (scalar BAC
in g/dL) & BiLSTM + attention, classification head (4-class softmax) \\
Output type & Scalar (float) & Ordinal class (0--3) + confidence \\
Actuator & Ignition relay (binary ALLOW/BLOCK) & Graduated cabin profile
(LED colour, fan \%, audio) \\
BLE packet size & 20 bytes & 12 bytes \\
Real-time latency & $<50$~ms on-device + 30~s emission cadence &
$<50$~ms on-device + 30~s emission cadence \\
Safety failure mode & Conservative: ignition blocked by default &
Conservative: WAITING state on any packet loss \\
Primary risk addressed & Alcohol-related accidents & Stress-related
accidents \\
\bottomrule
\end{tabular}
\end{table}

Three architectural decisions are shared between the two systems:
the on-device TensorFlow Lite inference pipeline, the Arduino-based
vehicle-side controller, and the conservative ``default deny''
failure mode in which the loss of a valid BLE update drives the
vehicle module into its safest state. Two are specific to the
present work: the classification (rather than regression) output
head, which better matches the graduated intervention philosophy
described in Chapter~\ref{ch:introduction}; and the shrinking of
the BLE packet from twenty to twelve bytes, reflecting that an
ordinal class occupies less space than a floating-point regression
output and that no explicit alert-level byte is needed because the
class already encodes the severity.

Chapter~\ref{ch:conclusion} returns to this comparison and sketches
a unified wearable safety platform that would emit both
characteristics side by side on the same BLE service.
